%\documentclass[11pt,a4paper,twoside]{article}

\usepackage[utf8]{inputenc}

\usepackage[T1,T2A]{fontenc}

\usepackage[english.ancient,russian.ancient]{babel}

\usepackage{graphicx}

\usepackage{array}

\usepackage[doi]{samgtu-bib}

\usepackage{samgtu}

\usepackage{samgtu-name}

\usepackage{mathtools}

\mathtoolsset{showonlyrefs}

\usepackage{desclist}

\usepackage{euscript}

\usepackage{mathrsfs} 

\usepackage{multicol}

\usepackage{mathrsfs}

\usepackage{cite}

\usepackage{cmap}

\usepackage{qrcode}

\allowdisplaybreaks[4]

\usepackage[nodayofweek]{datetime}

\usepackage[thinlines]{easybmat}



\renewcommand{\JournalYear}{2018}

\renewcommand{\JournalVolume}{22}

\renewcommand{\JournalNumber}{4}



\newdate{JournalDateSign}{29}{12}{2018}% Дата подписания Вестника в печать

\renewcommand{\JournalDateSign}{\displaydate{JournalDateSign}}

\newdate{JournalDatePrint}{16}{01}{2019}% Дата выхода Вестника из печати

\renewcommand{\JournalDatePrint}{\displaydate{JournalDatePrint}}

%\renewcommand{\JournalUPL}{16.56} % Усл. печ. л.

%\renewcommand{\JournalUIL}{16.53} % Уч.-изд. л.

\renewcommand{\JournalUPL}{15.24} % Усл. печ. л.

\renewcommand{\JournalUIL}{15.21} % Уч.-изд. л.

\renewcommand{\JournalRegNumber}{220/18} % Регистрационный номер

\renewcommand{\JournalOrderNumber}{2} % Номер заказа



%\renewcommand{\JournalRMonth}{Март} % Месяц выхода Вестника

%\renewcommand{\JournalEMonth}{March} % Месяц выхода Вестника

%\renewcommand{\JournalRMonth}{Июнь} % Месяц выхода Вестника

%\renewcommand{\JournalEMonth}{June} % Месяц выхода Вестника

%\renewcommand{\JournalRMonth}{Сентябрь} % Месяц выхода Вестника

%\renewcommand{\JournalEMonth}{September} % Месяц выхода Вестника

\renewcommand{\JournalRMonth}{Декабрь} % Месяц выхода Вестника

\renewcommand{\JournalEMonth}{December} % Месяц выхода Вестника



\newcommand{\totop}[1]{\raisebox{6pt}[1pt][2pt]{#1}}

\newcommand{\tottop}[1]{\raisebox{12pt}[1pt][2pt]{#1}} 

\newcommand{\totttop}[1]{\raisebox{18pt}[1pt][2pt]{#1}} 

\newcommand{\tottttop}[1]{\raisebox{24pt}[1pt][2pt]{#1}} 

\newcommand{\totttttop}[1]{\raisebox{30pt}[1pt][2pt]{#1}} 

\newcommand{\tobot}[1]{\raisebox{-6pt}[1pt][2pt]{#1}} 

\newcommand{\tobbot}[1]{\raisebox{-12pt}[1pt][2pt]{#1}} 

\newcommand{\tobbbot}[1]{\raisebox{-18pt}[1pt][2pt]{#1}}



\graphicspath{{./00PIC/}}

%

\vsgtu{1568} %registration number

\FirstPage{736}

\LastPage{751}

%

%\pdfcrop



%\draft

%

%

\ResearchArticle

%

%

%

%\begin{document}





%\hfill \qrcode[hyperlink,height=.8in]{http://doi.org/10.14498/vsgtu1568}

%\vspace{-.8in}



\renewcommand{\baselinestretch}{0.89}



\udc{517.958:532.51}

\msc{76F02, 76F45, 76M45, 76R05, 76U05}







\rutitle{Крупномасштабная слоистая стационарная конвекция~вязкой несжимаемой жидкости под~действием~касательных напряжений на~верхней~границе. Исследование полей температуры~и~давления}

\rutitleshort{Крупномасштабная слоистая стационарная конвекция вязкой несжимаемой жидкости\dots}



\entitle{A large-scale layered stationary convection of~a~incompressible viscous fluid under the action of shear stresses at the upper boundary. Temperature and presure field investigation}

\entitleshort{A large-scale layered stationary convection of a incompressible viscous fluid\dots}



\ruauthor{Н.~В.~Бурмашева$^{1,2}$, Е.~Ю.~Просвиряков$^2$}{Б\,у\,р\,м\,а\,ш\,е\,в\,а~Н.~В., П\,р\,о\,с\,в\,и\,р\,я\,к\,о\,в~Е.~Ю.}





\enauthor{N.~V.~Burmasheva$^{1,2}$, E.~Yu.~Prosviryakov$^2$}{B\,u\,r\,m\,a\,s\,h\,e\,v\,a~N.~V., P\,r\,o\,s\,v\,i\,r\,y\,a\,k\,o\,v~E.~Yu.}



\ruaffil{\ruaffid{7370}{Уральский федеральный университет \\ им.~первого Президента России Б.~Н.~Ельцина, \\ 

Россия, 620002, Екатеринбург, ул. Мира, 19}.

\and

\ruaffid{2665}{Институт машиноведения УрО РАН,\\

Россия, 620049, Екатеринбург, ул. Комсомольская, 34}.}

\enaffil{\enaffid{7370}{Ural Federal University named after the First President of Russia B.~N.~Yeltsin, \\ 19, Mira st., Ekaterinburg, 620002, Russian Federation}.

           \and

\enaffid{2665}{Institute of Engineering Science, Urals Branch, Russian Academy of Sciences,\\

34, Komsomolskaya st., Ekaterinburg, 620049, Russian Federation}.}





\ruauthorsdetails{2}{

\fullauthor{\rupid{52636}{Наталья Владимировна Бурмашева}}

\CAuthor

\orcid{0000-0003-4711-1894}

\regalia{кандидат технических наук} 

\position{доцент} 

\dept{институт математики и компьютерных наук, каф. прикладной математики и механики$^1$}

\position{научный сотрудник} 

\dept{сектор нелинейной вихревой гидродинамики$^2$}

\email{nat_burm@mail.ru}

\fullauthor{\rupid{41426}{Евгений Юрьевич Просвиряков}} 

\orcid{0000-0002-2349-7801}

\regalia{доктор физико-математических наук} 

\position{заведующий сектором} 

\dept{сектор нелинейной вихревой гидродинамики$^2$}

\email{evgen_pros@mail.ru}}







\enauthorsdetails{2}{

\fullauthor{\enpid{52636}{Natalya V. Burmasheva}} 

\CAuthor

\orcid{0000-0003-4711-1894}

\regalia{Cand. Techn. Sci.} 

\position{Associate Professor} 

\dept{Institute of Mathematics and Computer Science, Dept. of Applied Mathematics and Mechanics$^1$}

\position{Researcher} 

\dept{Sector of Nonlinear Vortex Hydrodynamics$^2$}

\email{nat_burm@mail.ru}

\fullauthor{\enpid{41426}{Evgeny Yu. Prosviryakov}} 

\orcid{0000-0002-2349-7801}

\regalia{Dr. Phys. \& Math. Sci.} 

\position{Head of Sector} 

\dept{Sector of Nonlinear Vortex Hydrodynamics$^2$}

\email[br]{evgen_pros@mail.ru}\vspace{-5mm}}



\rukeywords{система Обербека--Буссинеска, сдвиговое течение, конвекция, точное решение, полиномиальное решение, локализация корней, термоклин, расслоение полей}

\enkeywords{Oberbeck--Boussinesq system, shear flow, convection, exact solution, polynomial solution, root localization, thermocline, field stratification}



\ruabstract{Изучается новое точное решение переопределенной системы уравнений Обербека"--~Буссинеска, которое описывает стационарное сдвиговое течение вязкой несжимаемой жидкости в бесконечном слое.  Приведенное точное решение является обобщением класса Остроумова"--~Бириха для слоистого однонаправленного потока. В предложенном решении горизонтальные скорости зависят только от поперечной координаты $z$. Поле температуры и поле давление являются трехмерными. В отличие от решения Остроумова"--~Бириха, в представленном в статье решении горизонтальные градиенты температуры являются линейными функциями от координаты $z$. Такая структура точного решения позволяет найти нетривиальное решение уравнений  Обербека"--~Буссинеска посредством тождественного равенства нулю уравнения несжимаемости. Данное точное решение пригодно для  исследования крупномасштабных течений вязкой несжимаемой жидкости квазидвумерными уравнениями.  Конвективное движение жидкости обусловлено заданием касательных напряжений на свободной границе слоя. Неоднородные тепловые источники заданы на обеих границах. Давление в жидкости на верхней границе совпадает с атмосферным давлением. В статье основное внимание уделяется исследованию полей температуры и давления, которые описываются многочленами трех переменных. Детально обсуждаются особенностей распределения профилей температуры и давления, которые являются многочленами седьмой и восьмой степени соответственно. Для анализа свойств температуры и давления используются алгебраические методы для исследования числа корней на отрезке.  Показано, что фоновая температура и фоновое давление являются немонотонными функциями. Температурное поле расслаивается на зоны, которые формируют термоклин и тепловой пограничный слой около границ слоя жидкости. Исследование свойств поля давления показали, что оно расслаивается на одну, две или три зоны относительно отсчетного значения (атмосферного давления).}



\enabstract{In this paper a new exact solution of an overdetermined system of Ober\-beck--Boussinesq equations that describes a stationary shear flow of a viscous incompressible fluid in an infinite layer is under study. The given exact solution is a generalization of the Ostroumov--Birich class for a layered unidirectional flow. In the proposed solution, the horizontal velocities depend only on the transverse coordinate z. The temperature field and the pressure field are three-dimensional. In contradistinction to the Ostroumov--Birich solution, in the solution presented in the paper the horizontal temperature gradients are linear functions of the $z$ coordinate. This structure of the exact solution allows us to find a nontrivial solution of the Oberbeck--Boussinesq equations by means of the identity zero of the incompressibility equation. This exact solution is suitable for investigating large-scale flows of a viscous incompressible fluid by quasi-two-dimensional equations. Convective fluid motion is caused by the setting of tangential stresses on the free boundary of the layer. Inhomogeneous thermal sources are given on both boundaries. The pressure in the fluid at the upper boundary coincides with the atmospheric pressure. The paper focuses on the study of temperature and pressure fields, which are described by polynomials of three variables. The features of the distribution of the temperature and pressure profiles, which are polynomials of the seventh and eighth degree, respectively, are discussed in detail. To analyze the properties of temperature and pressure, algebraic methods are used to study the number of roots on a segment. It is shown that the background temperature and the background pressure are nonmonotonic functions. The temperature field is stratified into zones that form the thermocline and the thermal boundary layer near the boundaries of the fluid layer. Investigation of the properties of the pressure field showed that it is stratified into one, two or three zones relative to the reference value (atmospheric pressure).}



\newdate{burma}{17}{10}{2017}

\date{\displaydate{burma}}

\newdate{burmb}{15}{12}{2017}

\revisiondate{\displaydate{burmb}}

\newdate{burmc}{18}{12}{2017}

\accepteddate{\displaydate{burmc}}



\newdate{burme}{29}{12}{2017}

\renewcommand{\JournalDataOnline}{\displaydate{burme}}







\makerutitle





\Section[n]{Введение} 

Изучение конвективных течений вязкой несжимаемой жидкости  является одной из главных проблем при описании природных явлений и конструировании технических устройств. Нелинейность процессов перемешивания, протекающих в жидкости или газе, требует повышенного внимания к методам, позволяющим изучить распределение гидродинамических полей в расчетной области. Понятно, что для описания потоков несжимаемых сред необходимо иметь определенный запас точных решений уравнений движения.



При описании неизотермических течений вязкой несжимаемой жидкости в условиях нормальной гравитации наибольшее распространение получила конвекция, описываемая уравнениями Обербека"--~Буссинеска \cite{Joseph, Gershuni1989, Pukhnachev4}. 

Несмотря на приближенность уравнений Обербека"--~Буссинеска, в которых пренебрегают сжимаемостью вязкой жидкости, они очень сложно поддаются аналитическому исследованию.



Первый класс точных решений был предложен Г.~А.~Остроумовым \cite{Ostroumov} и Р.~В.~Бирихом \cite{Birikh}. Характерной особенностью семейства точных решений Остроумова"--~Бириха является равенство нулю вертикальной скорости и одной из горизонтальных скоростей. 

Таким образом, точное решение Остроумова"--~Бириха описывает однонаправленное конвективное течение $ (V_x; 0; 0)$ вязкой несжимаемой жидкости. Данное точное решение было получено повторно независимо от представителей пермской гидродинамической школы в~статьях \cite{Ortiz, Smith}.



К настоящему времени течение Остроумова"--~Бириха многократно обобщалось, использовалось при решении задач устойчивости и для описания гидрологических и технологических процессов \cite{Andreev, Andreev_bek, Andreev_step, Bek_Gon_Rez_Shef, Bir_Den_Kost1, Bir_Den_Kost2, Bir_Pukhn, Birikh2, Gon_Kab2, Gon_rez, Pukhnachev6}.

Другим направлением, обобщающим класс течений типа Остроумова"--~Бириха, является представление поля скоростей суперпозицией однонаправленных горизонтальных течений. 

В~этом случае слоистое движение жидкости $(V_x; V_y; 0 )$ имеет сдвиговый характер

\cite{Aristov_avto, Aristov_Pros3, Aristov_Pros5,  Aristov_Pros7, Aristov_Pros_Spev1, Aristov_Pros_Spev2, Burm_Pros_Samara, Burm_Pros_IOP, Burm_Pros_DREAM, Burm_Pros_AIP, Gorsh_Pros1, Gorsh_Pros2, Knyaz1, Sidorov1,  Schwartz1}.

Понятно, что указанное обобщение оказалось востребованным для решения задач геофизической гидродинамики. О новых свойствах конвективных течений вращающейся вязкой несжимаемой жидкости можно узнать из библиографических источников \cite{Arist2006, Aristov_Schwartz4, Schwartz2, Schwartz3}.



Точные решения типа Бириха"--~Остроумова и их модификации оказались эффективными для описания противотечений в~вязкой несжимаемой жидкости. 

Известно, что при крупномасштабных течениях несжимаемых сред противотечения в~жидкости могут иметь термическое происхождение, обусловленное нелинейным взаимодействием жидких частиц

\cite{Andreev_step, Bek_Gon_Rez_Shef, Bir_Den_Kost1, Bir_Den_Kost2, Bir_Pukhn, Birikh2, Gon_Kab2, Gon_rez, Aristov_Pros3, Aristov_Pros5, Aristov_Pros7, Aristov_Pros_Spev1, Aristov_Pros_Spev2, Burm_Pros_Samara, Burm_Pros_IOP,Burm_Pros_DREAM, Gorsh_Pros1, Gorsh_Pros2, Arist2006, Shtern}.

Для моделирования геофизических течений вязкой несжимаемой жидкости в~качестве граничных условий задаются касательные напряжения на границе атмосферы и океана. 

Очевидно, что тангенциальные силы могут быть вызваны различными факторами. Существенную роль играет и термокапиллярный эффект, влияние которого на противотечение было изучено в~статьях \cite{Aristov_Pros7, Aristov_Pros_Spev1, Aristov_Pros_Spev2, Gorsh_Pros1, Gorsh_Pros2}.



Обобщение результатов, анонсированных в \cite{Aristov_Pros7, Aristov_Pros_Spev1, Aristov_Pros_Spev2, Gorsh_Pros1, Gorsh_Pros2}, осуществлено в статье \cite{Burm_Pros_Samara}. Было показано, что поле скоростей может иметь две застойные точки, что является принципиальным отличием от слоистой конвекции Марангони. 

Учитывая нахождение новых физических эффектов при учете касательных напряжений при течении жидкости в бесконечном слое, в данной статье достаточно подробно изучаются поля температуры и давления.



\smallskip



\Section{Постановка задачи}

Система уравнений тепловой конвекции в приближении Буссинеска \cite{Gershuni1989}

в~случае стационарных сдвиговых течений $\left(V_x; V_y; 0\right)$  принимает вид

\begin{gather}%\label{eq:NS-x}

\displaystyle

V_x\frac{\partial V_x}{\partial x}+V_y\frac{\partial V_x}{\partial y}=-\frac{\partial P}{\partial x} +\nu\triangle V_x,\\

%\end{equation}

%\begin{equation}%\label{eq:NS-y}

\displaystyle

V_x\frac{\partial V_y}{\partial x}+V_y\frac{\partial V_y}{\partial y}=-\frac{\partial P}{\partial y} +\nu\triangle V_y,\\

%\end{equation}

%\begin{equation}

\label{eq:ur-conv}

\frac{\partial P}{\partial z}=g\beta T,\\

%\end{equation}

%\begin{equation}%\label{eq:teploprov}

\displaystyle

V_x\frac{\partial T}{\partial x}+V_y\frac{\partial T}{\partial y}=\chi\triangle T,\\

%\end{equation}

%\begin{equation}%\label{eq:nesgim}

\displaystyle

\frac{\partial V_x}{\partial x}+\frac{\partial V_y}{\partial y}=0.

\end{gather}

Здесь $P(x,y,z)$ "--- отклонение давления от гидростатического, деленное на постоянную среднюю плотность $\rho$ жидкости; $T(x,y,z)$ "--- отклонение от отсчетного значения температуры; $\nu,$ $\chi$ "--- коэффициенты кинематической вязкости и температуропроводности жидкости соответственно; $\triangle= \frac{\partial^2}{\partial x^2}+\frac{\partial^2}{\partial y^2}+\frac{\partial^2}{\partial z^2}$ "--- оператор Лапласа.



Система \eqref{eq:ur-conv}

%\eqref{eq:NS-x, eq:NS-y, eq::NS-z, eq:teploprov, eq:nesgim}

переопределена, поскольку из пяти уравнений в частных производных необходимо вычислить четыре функции. 

В~\cite{Arist2006, Burm_Pros_Samara, Burm_Pros_IOP, Burm_Pros_DREAM, Burm_Pros_AIP, Gorsh_Pros1, Gorsh_Pros2} показано, что эта система разрешима в классе

\begin{equation}\label{eq:klass-resh-skor}

V_x=u(z), \quad V_y=v(z).

\end{equation}

Поле скоростей \eqref{eq:klass-resh-skor} тождественно удовлетворяет уравнению несжимаемости системы \eqref{eq:ur-conv}, что позволяет исключить <<лишнее>> уравнение.



Поля давления и температуры в этом случае линейны по горизонтальным (продольным) координатам $x$ и $y$:

\begin{equation}\label{eq:klass-resh-polya}

T=T_0(z)+T_1(z)x+T_2(z)y,\quad P=P_0(z)+P_1(z)x+P_2(z)y.

\end{equation}



Подстановка разложения \eqref{eq:klass-resh-skor},  \eqref{eq:klass-resh-polya} в первое уравнение системы \eqref{eq:ur-conv} приводит к уравнению

$$\displaystyle

0=-\frac{\partial}{\partial x}

\bigl(P_0(z)+P_1(z)x+P_2(z)y \bigr)+\nu 

\Bigl(\frac{\partial^2}{\partial x^2}+\frac{\partial^2}{\partial y^2}+\frac{\partial^2}{\partial z^2}\Bigr) u(z).$$

Совершая очевидные преобразования и пользуясь методом неопределенных коэффициентов, получаем обыкновенное дифференциальное уравнение, описывающее баланс вязких сил и градиента давления 

$$\displaystyle

 \nu u^{\prime\prime}=P_1.$$

Здесь и далее штрих обозначает дифференцирование по координате $z$.



Аналогичным образом система уравнений тепловой конвекции  может быть приведена к системе обыкновенных дифференциальных уравнений вида

\begin{gather}%\label{eq:T12}

 T_1^{\prime\prime}=0,   \quad   T_2^{\prime\prime}=0,\\

%\end{equation}

%\begin{equation} %\label{eq:P12}

P_1^{\prime}=g\beta T_1,  \quad  P_2^{\prime}=g\beta T_2,  \\

%\end{equation}

%\begin{equation}

\label{eq:ur-conv-odu}%\label{eq:uv}

 \nu u^{\prime\prime}=P_1,  \quad   \nu v^{\prime\prime}=P_2,\\

%\end{equation}

%\begin{equation}%\label{eq:T0}

\chi T_0^{\prime\prime}=u T_1+v T_2,\\

%\end{equation}

%\begin{equation}%\label{eq:P0}

 P_0^{\prime}=g\beta T_0.

\end{gather}



Более подробно получение системы \eqref{eq:ur-conv-odu} изложено в \cite{Burm_Pros_Samara}.

Система уравнений \eqref{eq:ur-conv-odu} имеет общее решение, которое является полиномиальным \cite{Aristov_Pros5, Aristov_Pros7, Aristov_Pros_Spev1, Aristov_Pros_Spev2, Burm_Pros_Samara, Burm_Pros_IOP,Burm_Pros_DREAM,Burm_Pros_AIP,Gorsh_Pros1}.

%[ссылки на родственные статьи.].





Далее сформулируем граничные условия, определяющие конвективное течение вязкой несжимаемой жидкости в бесконечном слое. 

Полагаем, что нижняя граница бесконечного слоя жидкости является абсолютно твердой и неподвижной, а верхняя "--- свободной. 

Будем считать, что свободная поверхность является недеформируемой.



В качестве краевых условий рассмотрим следующие граничные возмущения полей скорости, температуры и давления \cite{Burm_Pros_Samara}:

\begin{itemize}

\item на нижней термоизолированной ($T(x,y,0)=0$) границе ($z=0$) слоя жидкости выполняется условие прилипания:

$$V_x (0)=V_y (0)=0;$$

\item на верхней границе ($z=h$) действует постоянное давление $${P(x,y,h)=S},$$ а~температура задается линейной функцией (нагрев или охлаждение)

$$T(x,y,h)=\vartheta+Cx+Dy;$$

\item кроме того, на свободной поверхности $z=h$ заданы касательные напряжения

$$\eta \displaystyle \frac{d V_x}{dz}=\xi_1, \quad \eta \displaystyle \frac{d V_y}{dz}=\xi_2.$$

\end{itemize}



Учитывая структуру класса точных решений \eqref{eq:klass-resh-skor},  \eqref{eq:klass-resh-polya}, сформулированные выше граничные условия записываются следующим образом:

\begin{gather}

u(0)=v(0)=0,\\

T_0(0)=0, \quad T_1(0)=0, \quad T_2(0)=0,\\

T_0(h)=\vartheta, \quad T_1(h)=C, \quad T_2(h)=D,\\

P_0(h)=S, \quad P_1(h)=0, \quad P_2(h)=0,\\

\displaystyle \eta\frac{du}{dz}(h)=\xi_1, \quad \eta\frac{dv}{dz}(h)=\xi_2.

\end{gather}

Без ограничения общности полагаем $S=0$. 

Другими словами, будем вести отсчет приведенного давления от уровня, задаваемого на верхней границе слоя жидкости.

%Подробно опиши каждую группу граничных условий и упомяни, что течение происходит в бесконечном слое жидкости.



\smallskip



\Section{Решение системы уравнений}\label{sec:solution}

Интегрирование системы \eqref{eq:ur-conv-odu} в~силу краевых условий дает точное решение \cite{Burm_Pros_Samara}:

%Я бы про интегрирование убрал, а выписал сразу точное решение для T_0, P_0, ссылаясь на статью в Самаре.

$$T_1=C\frac{z}{h}, \quad T_2=D\frac{z}{h},$$

\begin{multline}

T_0(z)=\displaystyle \frac{z \left(-h^3 (C \xi_1 + D \xi_2) +

    12 \chi \eta \vartheta + (C \xi_1 + D\xi_2) z^3 \right)}{12 \chi \eta h} +

\\

+\displaystyle \frac{C D g\beta z \left(-16 h^6 + 21 h^2 z^4 -

    5 z^6 \right)}{840 \chi h^2 \nu},

\end{multline}

$$P_1=\displaystyle g\beta C\frac{z^2-h^2}{2h}, \quad P_2=\displaystyle g\beta D\frac{z^2-h^2}{2h},$$

\begin{multline}

P_0(z)=\displaystyle \frac{\beta g \left[ \left(C \xi_1 + D\xi_2\right)\left(3 h^5-5 h^3z^2+2 z^5\right) - 60 \chi \eta \vartheta\left(h^2 - z^2\right)\right]}{120 \chi \eta h} +

\\

+\displaystyle \frac{ \beta^2 C D g^2 \left(41 h^8 - 64 h^6 z^2 + 28 h^2 z^6 - 5 z^8 \right)}{6720 \chi h^2 \nu}.

\end{multline}



Далее решение, описывающее температуру $T$ и давление $P,$ будем рассматривать в безразмерном виде:

\begin{gather}

T_1=Z, \quad T_2=K Z,

\\

T_0=\displaystyle Z(Z^3-1)\alpha+Z\gamma+\frac{\delta {\sf Pe}}{840}Z(-16+21Z^4-5Z^6),

\\

P_1=(Z^2-1 )\frac{\delta}{2}, \quad P_2=K\left(Z^2-1 \right)\frac{\delta}{2},

\\

P_0=\displaystyle \frac{\left(\xi_1+K\xi_2\right)h^4}{120\chi\eta l^2}\left(3-5Z^2+2Z^5\right)+\frac{\vartheta}{Cl}\frac{\delta}{2} \left(Z^2-1\right)+\hspace{3cm}\\

\hspace{5cm}+\displaystyle \frac{g\beta Dh^6}{6720\chi\nu l^2}\left(41-64Z^2+28Z^6-5Z^8\right).

\end{gather}

Здесь $K=D/C$, $\delta=h/l$, ${\sf Pe}={\sf Pr}\, {\sf Gr}$ "--- число Пекле, 

${\sf Pr}={\nu}/{\chi}$ "--- число Прандтля, ${\sf Gr}={Dhg\beta h^3}/{\nu^2}$ "--- число Грасгофа, 

$h$ "--- толщина слоя жидкости, $l$ "--- характерный размер вдоль горизонтальных координат. 

В~формулах для безразмерных точных решений сохранены обозначения такие же, как и для размерных полей температуры и~давления. 

Отметим, что коэффициенты 

$\frac{\delta {\sf Pe}}{840}$, 

$\frac{(\xi_1+K\xi_2) h^4}{ 120\chi\eta l^2}$, 

$\frac{\vartheta\delta}{2Cl}$, 

$\frac{g\beta Dh^6}{6720\chi\nu l^2}=\frac{\delta^2{\sf Gr\,Pr}}{6720}$, 

стоящие в разложении функции $T_0$ и $P_0$, являются безразмерными.

Процедура перехода к безразмерным гидродинамическим полям подробно изложена в работах \cite{Burm_Pros_Samara, Burm_Pros_IOP, Burm_Pros_DREAM, Burm_Pros_AIP}.



\smallskip





\Section{Исследование стратификации температурного поля} \label{sec:analiz-T-rassloen}

Поскольку поле температуры $T$ является полиномиальным, следовательно, непрерывным, то исследование его структуры можно свести к изучению нулей функции $T_0$. Многочлен $T_0$ в силу граничных условий можно представить в мультипликативном виде:

$$T_0(Z)=\displaystyle \frac{\delta {\sf Pe}}{840} Z f(Z),$$

где 

\begin{equation}

f(Z)=aZ^3-5Z^6+21Z^4-16+b, \quad a=\displaystyle \frac{840\alpha}{\delta {\sf Pe}}, \quad  b=\displaystyle \frac{840\gamma}{\delta {\sf Pe}}-a.

\label{sec:f(Z)}

\end{equation}

Заметим, что $f(0)=b-16, f(1)=a+b$.



Поскольку введенная выше переменная $Z$ определена на отрезке $[0, 1]$, для определения количества нулей многочлена $f(Z)$ можно воспользоваться теоремой Декарта \cite{Dekart}. Согласно этой теореме, число положительных корней многочлена с вещественными коэффициентами равно количеству перемен знаков в ряду его коэффициентов или на четное число меньше этого количества.



Представим многочлен $f(Z)$ в стандартном виде

$$f(Z)=-5Z^6+21Z^4+aZ^3+(b-16)$$

и посчитаем число перемен знаков в получившемся ряду коэффициентов $\{-5; 21; a; (b-16)\}$, которое зависит от значений параметров $a$ и $b$. 

Несложно убедиться, что перемен знаков (следовательно, и число положительных корней многочлена $f(Z)$) может быть равно $1$, $2$ или $3$. И, следовательно, число положительных корней многочлена $f(Z)$ может быть равно $0$, $1$, $2$ или $3$. Таким образом, на положительной полуоси $( Z>0 )$ у функции $T_0$ может быть не более трех нулевых значений. 

Данное условие является необходимым для существования стратификации температурного поля относительно поперечной координаты $Z$.



Для того чтобы получить достаточное условие для оценки числа нулей многочлена $f(Z)$ на интервале $(0, 1)$, исследуем функцию $f(Z)$ на наличие экстремумов. 

Очевидно, что внутри интервала $(0, 1)$ их не более трех. Известно, что число экстремумов функции одной переменной определяется числом нулей ее производной:

$$f^{\prime}(Z)=-30Z^5+84Z^3+3aZ^2=Z^2\left(-30Z^3+84Z+3a \right).$$









\noindent

Здесь в выражении для $f^{\prime}(Z)$ в круглых скобках   стоит полином третьей степени, следовательно, число его действительных нулей не превосходит трех. При исследовании многочлена на экстремум в зависимости от значений параметров возможны различные случаи. Если  у функции $f$ нет экстремумов на интервале $(0,1)$ (строго монотонна), то все определяется значениями функции на концах отрезка~$[0,1]$:

\begin{itemize}

              \item если  $f(0)f(1)\geq 0$, то расслоений не возникает;

              \item если $f(0)f(1)<0$, то имеет место одна точка расслоения.

\end{itemize}







Аналогичным образом могут быть проанализированы ситуации, когда у функции $f(Z)$ есть один, два или три экстремума.

Как отмечено выше, число экстремумов функции определяется числом нулей ее производной:

$$f^{\prime}(Z)=-30Z^5+84Z^3+3aZ^2=3Z^2 (-10Z^3+28Z+a)=0;$$

$$Z=0 \quad  \text{или} \quad  10Z^3-28Z-a=0.$$



Если при некоторых условиях на $a$ последнее соотношение выполняется хотя бы в одной точке на интервале $(0,1)$, то у функции $f$ на интервале $(0,1)$ есть экстремумы (столько, сколько точек, в которых выполняется условие $10Z^3-28Z-a=0$), иначе функция $f$ на этом интервале строго монотонна.





Рассмотрим семейство функций $g_a(Z)=10Z^3-28Z-a$ и обозначим через $g(Z)$ функцию этого семейства при $a=0$. 

%Графики этих функций представлены на рис.~\ref{pic:ga}. 

%\begin{figure}[h!]

%\begin{center}

%\begin{minipage}[h!]{0.49\linewidth}

%  \centering

%  % \psfrag{Z}{$Z$}

%  % % \psfrag{g}{$g_a(Z)$}

%  \includegraphics[width=\textwidth]{ga(z)}

%\end{minipage}

%%\\

%\hfill

%\begin{minipage}[h!]{0.49\linewidth}

%  \centering

%  \caption{Семейство функций $g_a(Z)$ при разных $a$}\label{pic:ga}

%  \smallskip \footnotesize

%

%  [Figure~\ref{pic:ga}. The family of functions $g_a(Z)$ for different $a$]

%\end{minipage}

%\end{center}

%\end{figure}

%

Все функции семейства $g_a(Z)$  имеют  на интервале $(0,1)$ единственный экстремум в точке $Z_{\rm extr}=\sqrt{14/15}$.

Зная число экстремумов функций семейства $g_a(Z)$, можно, согласно приведенным выше рассуждениям, определить число нулей функций $g_a(Z)$ на интервале  $(0,1)$. А каждый такой нуль есть суть экстремум функции $f(Z)$ на этом интервале:

%($g_a^{\prime}(z)=g^{\prime}(z)=30z^2-28$).

%Таким образом,

%\begin{center}

%\begin{table}[h!]

%\begin{tabular}{rl}

\begin{itemize}

   \item  если $a < g(Z_{\rm extr})=-\frac{56}{3}\sqrt{\frac{14}{15}}$, то $g_a(Z)\neq0$ на интервале $(0,1)$, а функция $f(Z)$ "--- строго монотонна на интервале $(0,1)$; 

    \item  если $a = g(Z_{\rm extr})$,  то $g_a(Z)$ обращается в нуль в единственной точке, а функция $f(Z)$ в этой точке имеет экстремум; 

    \item  если $g(Z_{\rm extr})<a<g(1)=-18$, то $g_a(Z)$ имеет 2 нуля на интервале $(0, 1)$, а $f(Z)$ в~этих двух точках имеет экстремумы; 

    \item  если $g(1)\leq a\leq g(0)=0$, то $g_a(Z)$ обращается в нуль в единственной точке, а $f(Z)$ в этой точке имеет экстремум; 

    \item  если $g(0)<a$, то $g_a(Z)\neq 0$ на интервале $(0, 1)$, а функция $f(Z)$ "--- строго монотонна  на интервале $(0, 1)$.

%\end{tabular}

%\caption{___________}

%\label{__________}

%\end{table}

%\end{center}

\end{itemize}



Теперь, зная число экстремумов у функции $f(Z)$, можно выписать условия, при которых эта функция имеет нули на интервале $(0, 1)$. Эти нули, в свою очередь, определяют точки расслоения фоновой температуры. В качестве наглядной иллюстрации на рис.~\ref{pic:T0-all} приведены ситуации, соответствующие различным значениям параметров $a$ и $b$, при которых функция температуры допускает различное число расслоений.

%\ref{pic:T0-0}--\ref{pic:T0-2}),











\smallskip





\Section{Исследование расслоений для поля давления}

Исследуем теперь расслоение поля давления. Поле давления $P$ перепишем в виде

$$\displaystyle P(Z)= \frac{\delta^2 {\sf Pe}}{6720}(Z-1)q(Z),$$

где

\begin{multline}

q(Z)=-5Z^7-5Z^6+23Z^5+23Z^4+23Z^2-41Z-41

+

\\

+\phi\left(2Z^4+2Z^3+2Z^2-3Z-3\right)+\psi\left(Z+1\right)=

\\

=-5Z^7-5Z^6+23Z^5+(23+2\phi)Z^4+(23+2\phi)Z^3+

\\

+(23+2\phi)Z^2+(-41+\psi-3\phi)Z+(-41+\psi-3\phi),

\end{multline}

$$\displaystyle \phi=\frac{672\alpha}{\sf Pe}, \quad \psi=\frac{3360\gamma}{\delta {\sf Pe}}.$$





Согласно теореме Декарта \cite{Dekart}, число нулей полинома $q(Z)$ не превосходит трех, так как именно столько перемен знаков может быть в этом ряду при различных сочетаниях значений выражений $(23+2\phi)$ и $(-41+\psi-3\phi)$.



На рис.~\ref{pic:GMT} представлено геометрическое место точек, в которых функция $q(Z)$ обращается в нуль. Получившаяся поверхность седловая.

На рис.~\ref{pic:q123} представлено поведение функций $q_1(Z)= -5 Z^7 - 5 Z^6 + 23 Z^5$, $q_2(Z) = Z^4 + Z^3 + Z^2$, $q_3(Z) = Z + 1$, входящих в состав функции $q(Z)$:

$$

q(Z)=q_1(Z)+(23+2\phi) q_2(Z)+ (-41+\psi-3\phi)q_3(Z).

$$

 Отметим, что все три функции положительны и строго возрастают на отрезке $[0,1]$. 

 Поэтому выражения $(3+2\phi)$ и $(-41+\psi-3\phi)$ не могут быть одновременно положительными, иначе у функции $q(Z)$ на интересующем отрезке не будет нулей (она будет строго положительна).



\begin{figure}[p]



\centering \footnotesize



\includegraphics[scale=0.4]{burm-pict2}



\caption{Поведение функции $T_0(Z)$: нет расслоений ($a=-20, b=30$; линия {\sl 1}\/), одно расслоение ($a=-20, b=17$; линия {\sl 2}\/), два расслоения ($a=-14, b=16.5$; линия {\sl 3}\/);  

$a$ и $b$ см.~в~(\ref{sec:f(Z)}) \label{pic:T0-all} }



\smallskip



[Figure~\ref{pic:T0-all}.  The behavior of the function $T_0(Z)$:  no stratifications ($a=-20, b=30$; line~{\sl 1}\/), one stratification ($a=-20, b=17$; line~{\sl 2}\/), and two stratifications ($a=-14, b=16.5$; line~{\sl 3}\/); $a$~and~$b$~see in Eq.~\eqref{sec:f(Z)}]





\vspace{1cm}



\includegraphics[scale=0.4]{burm-pict3}

\caption{Геометрическое место точек, удовлетворяющих условию $q(Z)=0$ \label{pic:GMT} }



\smallskip

[Figure~\ref{pic:GMT}.   The geometric locus of points satisfying condition $q(Z)=0$]





\end{figure}







%

%\begin{figure}[h!]

%\begin{center}

%\begin{minipage}[h!]{0.49\linewidth}

%    \centering

%    % % \psfrag{z}{$Z$}

%    % % \psfrag{a}{\vspace{-1.5cm}\rotatebox{60}{\hspace{-0.5cm} $23+2\phi$}}

%    % % \psfrag{b}{\vspace{-1cm}\rotatebox{80}{$-41+\psi+3\phi$}}

%    \includegraphics[width=0.8\textwidth]{P0[z,a,b]==0}

%    \caption{Геометрическое место точек, удовлетворяющих условию $q(Z)=0$.}\label{pic:GMT}

%    \smallskip \footnotesize

%

%    [Figure~\ref{pic:GMT}.   The geometric locus of points satisfying condition $q(Z)=0$.]

%\end{minipage}

%\begin{minipage}[h!]{0.49\linewidth}

%    \centering

%    % % \psfrag{z}{$Z$}

%    %% % \psfrag{q1q2q3}{\hspace{0.5cm} $q_1(z), q_2(z), q_3(z)$}

%    % % \psfrag{q1q2q3}{}

%    \includegraphics[width=\textwidth]{q1q2q3}

%    \caption{Графики функций $q_1(Z)= -5 Z^7 - 5 Z^6 + 23 Z^5$ (жирная линия), $q_2(Z) = Z^4 + Z^3 + Z^2$, $q_3(Z) = Z + 1$ (пунктирная линия)}\label{pic:q123}

%    \smallskip \footnotesize

%

%    [Figure~\ref{pic:q123}. The graphs of functions $q_1(Z)= -5 Z^7 - 5 Z^6 + 23 Z^5$ (thick line), $q_2(Z) = Z^4 + Z^3 + Z^2$, $q_3(Z) = Z + 1$ (dashed line)]

%\end{minipage}

%\end{center}

%\end{figure}





На рис.~\ref{pic:P0-all} %\ref{pic:P0-0}--\ref{pic:P0-2}),

приведены различные случаи распределения фонового давления. Каждая из приведенный кривых иллюстрирует локализацию давления на верхней и нижней границах. Давление заметно отличается от гидростатического, что говорит о стратификации поля давления, обусловленной нелинейными механизмами переноса импульса в жидкости при конвективном перемешивании.



%\begin{figure}[h!]

%\noindent

%\begin{minipage}[h!]{0.5\linewidth}

%  \centering  \scriptsize \sl

%  % % \psfrag{z}{$Z$}

%  % % \psfrag{P}{$P_0(Z)$}

%  \includegraphics[width=0.98\textwidth]{P0-0,1,2}

%\end{minipage}

%\hfill

%\begin{minipage}[h!]{0.49\linewidth}

%  \caption{Поведение функции $P_0(Z)$: нет расслоений, одно расслоение (пунктирная линия), два расслоения (жирная линия}\label{pic:P0-all}

%  \smallskip \footnotesize

%

%  [Figure~\ref{pic:P0-all}.  The behavior of the function $P_0(Z)$: no stratifications, one stratification (dashed line), two stratifications (thick line)]

%\end{minipage}

%\end{figure}

%



\Section[n]{Заключение}

В статье получено новое точное решение уравнений Обербе\-ка"--~Буссинеска, описывающее конвективное сдвиговое течение вязкой несжимаемой жидкости. Конвекция в жидкости индуцируется заданием теплового источника и касательных напряжений на верхней (свободной) границе бесконечного слоя вязкой не\-сжимаемой жидкости. Приведенное в статье точное решение позволяет автоматически разрешить переопределенную систему уравнений переноса импульсов в теплопроводящей несжимаемой жидкости. Анализ полей температуры и давления, которые являются полиномиальными, показывает, что в вязкой несжимаемой жидкости происходит стратификация этих полей. Другими словами, характер распределения фоновой температуры не является монотонным, она может принимать одно или два нулевых значения, принятых за отсчетную конфигурацию. При определенных значениях граничных условий при течении жидкости могут наблюдаться локализованные решения, описывающие термоклин и пикноклин в несжимаемой жидкости.











\begin{figure}[p]



{

\centering \footnotesize



\includegraphics[scale=0.4]{burm-pict4}



\caption{Графики функций $q_1(Z)= -5 Z^7 - 5 Z^6 + 23 Z^5$ (линия~{\sl 1}\/), $q_2(Z) = Z^4 + Z^3 + Z^2$ (линия~{\sl 2}\/), $q_3(Z) = Z + 1$ (линия~{\sl 3}\/) \label{pic:q123} }

    \smallskip 



[Figure~\ref{pic:q123}. The graphs of functions $q_1(Z)= -5 Z^7 - 5 Z^6 + 23 Z^5$ (line~{\sl 1}\/), $q_2(Z) = Z^4 + Z^3 + Z^2$ (line~{\sl 2}\/), and $q_3(Z) = Z + 1$ (line~{\sl 3}\/)]





\vspace{0.5cm}



\includegraphics[scale=0.4]{burm-pict5}





\caption{Поведение функции $P_0(Z)$: нет расслоений (линия~{\sl 1}\/), одно расслоение (линия~{\sl 2}\/), два расслоения (линия~{\sl 3}\/) \label{pic:P0-all} }

  \smallskip 

  

[Figure~\ref{pic:P0-all}.  The behavior of the function $P_0(Z)$: no stratifications (line~{\sl 1}\/), one stratification (line~{\sl 2}\/), and two stratifications (line~{\sl 3}\/)]

%\end{minipage}

%\end{figure}

%



\bigskip



}



\AdditionalContent{Конкурирующие интересы}{Мы заявляем, что у нас нет конфликта интересов в отношении авторства и публикации этой статьи.}

\AdditionalContent{Авторский вклад и ответственность}{Мы  несем  полную  ответственность  за  предоставление  окончательной рукописи в печать. Каждый из нас одобрил окончательную  версию рукописи.}





\AdditionalContent{Финансирование}{Работа выполнена при поддержке фонда содействия развитию малых форм предприятий в научно-технической сфере (программа УМНИК), договор №~12281ГУ/2017.}





\end{figure}



\pagebreak





\begin{thebibliography}{99} 



\Bibitem{Joseph}

\by Joseph~D.~D.

\book Stability of fluid motions

\yr 1976

\publ Springer–Verlag

\publaddr Berlin, Heidelberg, New York

\totalpages 282





\RBibitem{Gershuni1989}

\by Гершуни~Г.~З., Жуховицкий~Е.~М., Непомнящий~А.~А.

\book Устойчивость конвективных течений

\yr 1989

\publ Наука

\publaddr М.

\totalpages 320

%\misctransl

%\lang In Russian

%\by Gershuni~G.~z., Zhuhovitskiy~E.~M., Nepomnyaschiy~A.~A.

%\book Ustoichivost' konvectivnikh techeniy \rm [Stability of Convective Flows]

%\yr 1989

%\publ Nauka

%\publaddr Moscow

%\totalpages 320



\RBibitem{Pukhnachev4}

\by Пухначёв~В.~В.

\paper Иерархия моделей в теории конвекции

\inbook Краевые задачи математической физики и смежные вопросы теории функций\rm.~32

\serial Зап. научн. сем. ПОМИ

\yr 2002

\vol 288

\pages 152--177

\publ ПОМИ

\publaddr СПб.

\mathnet{http://mi.mathnet.ru/znsl1588}

\mathscinet{http://www.ams.org/mathscinet-getitem?mr=1923549}

\zmath{https://zbmath.org/?q=an:1079.76644}

%\transl

%\jour J. Math. Sci. (N. Y.)

%\yr 2004

%\vol 123

%\issue 6

%\pages 4607--4620

%\crossref{10.1023/B:JOTH.0000041478.45024.64}



\RBibitem{Ostroumov}

\by Остроумов~Г.~А.

\book Свободная конвекция в условиях внутренней задачи

\yr 1952

\publ Гостехиздат

\publaddr М.

\totalpages 256

%\misctransl

%\lang In Russian

%\by Ostroumov~G.~A.

%\book Svobodnaya konvekciya v usloviyah vnutrenney zadachi \rm [Free convection under internal conditions]

%\yr 1952

%\publ Gostekhizdat

%\publaddr Moscow

%\totalpages 256



\RBibitem{Birikh}

\by Бирих~Р.~В.

\paper О термокапиллярной конвекции в горизонтальном слое жидкости

\jour ПМТФ

\yr 1966

\vol 7

\issue 3

\pages 69--72

%\misctransl

%\lang In Russian

%\by Birikh~R.~V.

%\paper Thermocapillary convection in a horizontal layer of liquid

%\jour Prikladnaya mehanika i thenicheskaya fizika \rm [J. Appl. Mech. Tech.Phys]

%\yr 1966

%\vol 7

%\issue 3

%\pages 69--72





\Bibitem{Smith}

\by Smith~M.~K., Davis~S.~H.

\paper Instabilities of dynamic thermocapillary liquid layers. Pt. 1. Convective instabilities

\jour J. Fluid Mech.

\yr 1983

\vol 132

\pages 119--144

\crossref{10.1017/S0022112083001512}



\Bibitem{Ortiz}

\by Ortiz-P\'{e}rez~A.~S., D\'{a}valos-Orozco~L.~A.

\paper Convection in a horizontal fluid layer under an inclined temperature gradient

\jour Phys. Fluids

\yr 2011

\vol 23

\issue 8

\pages 084107–084111

\crossref{10.1063/1.3626009}





\RBibitem{Andreev}

\by Андреев~В.~К.

\book Решения Бириха уравнений конвекции и некоторые его обобщения

\yr 2010

\publ ИВМ СО РАН

\publaddr Красноярск

\totalpages 68

%\misctransl

%\lang In Russian

%\by Andreev~V.K.

%\book Resheniya Biriha uravneniy konvekcii i nekotorie ego obobscheniya \rm [Birikh Solutions to Convection Equations and Some of its Extensions]

%\yr 2010

%\publ IBM SO RAN Publ.

%\publaddr Krasnoyarsk

%\totalpages 68



\RBibitem{Andreev_bek}

\by Андреев~В.~К., Бекежанова~В.~Б.

\paper Устойчивость неизотермических жидкостей (обзор)

\jour ПМТФ

\yr 2013

\vol 54

\issue 2

\pages 3--20

%\misctransl

%\by Andreev~V.~K., Bekezhanova~V.~B.

%\paper Stability of non-isothermal fluids (Review)

%\jour Journal of Applied Mechanics and Technical Physics

%\yr 2013

%\vol 54

%\issue 2

%\pages 171--184

%\crossref{10.1134/S0021894413020016}





\RBibitem{Andreev_step}

\by Андреев~В.~К., Степанова~И.~В.

\paper Однонаправленные течения бинарных смесей в моделе Обербека--Буссинеска

\jour  Изв. РАН. МЖГ

\yr 2016

\vol 51

\issue 2

\pages 13--24

\elib{https://elibrary.ru/item.asp?id=25934464}

%\misctransl

%\by Andreev~V.~K., Stepanova~I.V.

%\paper Unidirectional flows of binary mixtures within the framework of the Oberbeck–Boussinesq model

%\jour Fluid Dyn.

%\yr 2016

%\vol 51

%\issue 2

%\pages 136--147

%\crossref{10.1134/S0015462816020022}



\RBibitem{Bek_Gon_Rez_Shef}

\by Бекежанова~В.~Б., Гончарова~О.~Н., Резанова~Е.~В., Шефер~И.~А.

\paper Устойчивость двухслойных течений жидкости с испарением на границе раздела

\jour Изв. РАН. МЖГ

\yr 2017

\issue 2

\pages 23--35

\elib{https://elibrary.ru/item.asp?id=29008227}

\crossref{10.7868/S0568528117020062}

%\misctransl

%%\lang In Russian

%\by Bekezhanova~V.~B., Goncharova~O.~N., Rezanova~E.~V., Shefer~I.~A.

%\paper Stability of two-layer fluid flows with evaporation at the interface

%\jour Fluid Dyn.

%\yr 2017

%\vol 52

%\issue 2

%\pages 23--35

%\crossref{10.1134/S001546281702003X }

%%\zmath{https://zbmath.org/?q=an:1306.74011}

%%\scopus{http://www.scopus.com/record/display.url?origin=inward&eid=2-s2.0-84904339295}



\RBibitem{Bir_Den_Kost1}

\by Бирих~Р.~В., Денисова~М.~О.,  Костарев~К.~Г.

\paper Возникновение конвекции Марангони, вызванной локальным внесением поверхностно активного вещества

\jour Изв. РАН. МЖГ

\yr 2011

\issue 6

\pages 56--68

\elib{https://elibrary.ru/item.asp?id=17027130}

%\misctransl

%%\lang In Russian

%\by Birih~R.~V., Denisova~M.~O., Kostarev~K.G.

%\paper The development of marangoni convection induced by local addition of a surfactant

%\jour Fluid Dyn.

%\yr 2011

%\vol 46

%\issue 6

%\pages 890--900

%\crossref{10.1134/S0015462811060068 }





\RBibitem{Bir_Den_Kost2}

\by Бирих~Р.~В., Денисова~М.~О.,  Костарев~К.~Г.

\paper Развитие концентрационно-капиллярной конвекции на межфазной поверхности

\jour Изв. РАН. МЖГ

\yr 2015

\issue 3

\pages 56--67

\elib{https://elibrary.ru/item.asp?id=23774917}

%\misctransl

%%\lang In Russian

%\by Birih~R.~V., Denisova~M.~O., Kostarev~K.G.

%\paper Development of concentration-capillary convection on an interfacial surface

%\jour 	Fluid Dyn.

%\yr 2015

%\vol 50

%\issue 3

%\pages 361--370

%\crossref{10.1134/S0015462815030060 }





\RBibitem{Bir_Pukhn}

\by Бирих~Р.~В., Пухначёв~В.~В.

\paper Осевое конвективное течение во вращающейся трубе с продольным градиентом температуры

\jour Докл. РАН

\yr 2011

\vol 436

\issue 3

\pages 323--327

%\misctransl

%\lang In Russian

%\by Birih~R.~V., Puhnachev~V.~V.

%\paper An axial convective flow in a rotating tube with a longitudinal tem-perature gradient

%\jour Dokl. Phys.

%\yr 2011

%\vol 56

%\issue 1

%\pages 47--52

%\crossref{10.1134/S1028335811010095}





\RBibitem{Birikh2}

\by Бирих~Р.~В., Пухначёв~В.~В.

\paper Конвективное течение в горизонтальном канале с неньютоновской реологией поверхности при нестационарном продольном градиенте температуры

\jour Изв. РАН. МЖГ

\yr 2015

\issue 1

\pages 192--198

\elib{https://elibrary.ru/item.asp?id=23284010}

%\misctransl

%%\lang In Russian

%\by Birih~R.~V., Puhnachev~V.~V.

%\paper Convective flow in a horizontal channel with non-newtonian surface rheology under time-dependent longitudinal temperature gradient

%\jour Fluid Dyn.

%\yr 2015

%\vol 50

%\issue 1

%\pages 173--179

%\crossref{10.1134/S0015462815010172}





\RBibitem{Gon_Kab2}

\by  Гончарова~О.~Н.,  Кабов~О.~А.

\paper Гравитационно-термокапиллярная конвекция в горизонтальном слое при спутном потоке газа

\jour Докл. РАН

\yr 2009

\vol 426

\issue 2

\pages 183--188

\mathnet{http://mi.mathnet.ru/dan72}

\mathscinet{http://www.ams.org/mathscinet-getitem?mr=2536524}

%\misctransl

%\lang In Russian

%\by Goncharova~O.N., Kabov~O.A.

%\paper Gravitational-thermocapillary convection of fluid in the horizontal layer in co-current gas flow

%\jour Dokl. Phys.

%\yr 2009

%\vol 54

%\issue 5

%\pages 242--247

%\crossref{10.1134/S1028335809050061}



\RBibitem{Gon_rez}

\by Гончарова~О.~Н.,  Резанова~Е.~В.

\paper Пример точного решения стационарной задачи о двухслойных течениях с испарением на границе раздела

\jour ПМТФ

\yr 2014

\issue 2

\pages 68--79

%\misctransl

%%\lang In Russian

%\by Goncharova~O.~N., Rezanova~E.~V.

%\paper 	Example of an exact solution of the stationary problem of two-layer flows with evaporation at the interface

%\jour J. Appl. Mech. Tech.Phys

%\yr 2014

%\vol 55

%\issue 2

%\pages 247--257

%\crossref{10.1134/S0021894414020072}





\RBibitem{Pukhnachev6}

\by Пухначёв~В.В.

\paper Нестационарные аналоги решения Бириха

\jour Известия АлтГУ

\yr 2011

\issue 1--2

\pages 62--69

\elib{https://elibrary.ru/item.asp?id=17103846}

%\misctransl

%\lang In Russian

%\by Puhnachev~V.~V.

%\paper Non-stationary analogues of the Birikh solution

%\jour  Izvestiya AltGU \rm [Izvestiya AltGU]

%\yr 2011

%%\vol 54

%\issue 1--2

%\pages 62--69

%%\crossref{10.2478/s13540-014-0193-1}

%%\zmath{https://zbmath.org/?q=an:1306.74011}

%%\scopus{http://www.scopus.com/record/display.url?origin=inward&eid=2-s2.0-84904339295}



\RBibitem{Aristov_avto}

\by Аристов~С.~Н.

\thesis Вихревые течения в тонких слоях жидкости 

\thesisinfo Автореф. дис. \dots\ докт. физ.-мат. наук: 01.02.05

\yr 1990

\publ ИАПУ

\publaddr Владивосток

\totalpages 32





\RBibitem{Aristov_Pros3}

\by  Аристов~С.~Н.,  Просвиряков~Е.~Ю.

\paper Неоднородное конвективное течение Куэтта

\jour Известия Российской академии наук. Механика жидкости и газа

\yr 2016

\issue 5

\pages 3--9

\crossref{10.7868/S0568528116050030}

\elib{https://elibrary.ru/item.asp?id=26932326}

%\misctransl

%%\lang In Russian

%\by Aristov~S.N., Prosviryakov~E.~Yu.

%\paper Nonuniform convective Couette flow

%\jour Fluid Dyn.

%\yr 2016

%\vol 51

%\issue 5

%\pages 581--587

%\crossref{10.1134/S001546281605001X}



\RBibitem{Aristov_Pros5}

\by  Аристов~С.~Н.,  Просвиряков~Е.~Ю.

\paper Нестационарные слоистые течения завихренной жидкости

\jour Изв. РАН. МЖГ

\yr 2016

\issue 2

\pages 25--31

\elib{https://elibrary.ru/item.asp?id=25934465}

%%\zmath{https://?????}

%\misctransl

%%\lang In Russian

%\by Aristov~S.N., Prosviryakov~E.~Yu.

%\paper Unsteady layered vortical fluid flows

%\jour Fluid Dyn.

%\yr 2016

%\vol 51

%\issue 2

%\pages 148--154

%\crossref{10.1134/S0015462816020034}





\RBibitem{Aristov_Pros7}

\by  Аристов~С.~Н.,  Просвиряков~Е.~Ю.

\paper О слоистых течениях плоской свободной конвекции

\jour Нелинейная динам.

\yr 2013

\vol 9

\issue 4

\pages 651--657

\mathnet{http://mi.mathnet.ru/nd412}

\crossref{10.20537/nd1304004}

%\misctransl

%\lang In Russian

%\by Aristov~S.~N., Prosviryakov~E.~Yu.

%\paper On laminar flows of planar free convection

%\jour  Nelineynaya dinamika \rm [Nonlinear Dynamics]

%\yr 2013

%\vol 9

%\issue 4

%\pages 651--657

%%\crossref{10.20537/nd1304004}

%%\zmath{https://zbmath.org/?q=an:1306.74011}

%%\scopus{http://www.scopus.com/record/display.url?origin=inward&eid=2-s2.0-84904339295}





\RBibitem{Aristov_Pros_Spev1}

\by  Аристов~С.~Н.,  Просвиряков~Е.~Ю.,  Спевак~Л.~Ф.

\paper Нестационарная конвекция Бенара–Марангони слоистых течений вязкой несжимаемой жидкости

\jour Теоретические основы химической технологии

\yr 2016

\vol 50

\issue 2

\pages 137--146

\crossref{10.7868/S0040357116020019}

%\misctransl

%%\lang In Russian

%\by Aristov~S.~N., Prosviryakov~E.~Yu., Spevak~L.F.

%\paper Unsteady-state Bénard–Marangoni convection in layered viscous incompressible flows

%\jour  Teoreticheskie osnovi himicheskoy thenologii \rm [Theor. Found. Chem. Eng.]

%\yr 2016

%\vol 50

%\issue 2

%\pages 137--146

%\crossref{10.1134/S0040579516020019}

%%\zmath{https://zbmath.org/?q=an:1306.74011}







\RBibitem{Aristov_Pros_Spev2}

\by  Аристов~С.~Н.,  Просвиряков~Е.~Ю.,  Спевак~Л.~Ф.

\paper Нестационарная слоистая тепловая и концентрационная конвекция Марангони вязкой несжимаемой жидкости

\jour Вычислительная механика сплошных сред

\yr 2015

\vol 8

\issue 4

\pages 445--456

\crossref{10.7242/1999-6691/2015.8.4.38}

%\misctransl

%\lang In Russian

%\by Aristov~S.~N., Prosviryakov~E.~Yu., Spevak~L.F.

%\paper Nonstationary laminar thermal and solutal Marangoni convection of a viscous fluid

%\jour  Vichislitel'naya mehanika sploshnih sred \rm [Institute of Continuous Media Mechanics]

%\yr 2015

%\vol 8

%\issue 4

%\pages 445--456

%\crossref{10.7242/1999-6691/2015.8.4.38}







\RBibitem{Burm_Pros_Samara}

\by Бурмашева~Н.~В., Просвиряков~Е.~Ю.

\paper Крупномасштабная слоистая стационарная конвекция~вязкой несжимаемой жидкости

под~действием~касательных напряжений на~верхней~границе. Исследование поля скоростей

\jour Вестн. Сам. гос. техн. ун-та. Сер. Физ.-мат. науки

\yr 2017

\vol 21

\issue 1

\pages 180--196 

\mathnet{http://mi.mathnet.ru/vsgtu1527}

\crossref{10.14498/vsgtu1527}

\elib{http://elibrary.ru/item.asp?id=29245104}

\isi{000417031400010}

%\misctransl

%\lang In Russian

%\by Burmasheva~N.~V.,  Prosviryakov~E.~Yu.

%\paper A large-scale layered stationary convection

%of an incompressible viscous fluid under the action of shear

%stresses at the upper boundary. Velocity field investigation

%\jour Vestn. Samar. Gos. Tekh. Univ., Ser. Fiz. Mat. Nauki \rm [J. Samara State Tech. Univ., Ser. Phys. \& Math. Sci.]

%\yr 2017

%\vol 21

%\issue 1

%\pages 180--196

%\mathnet{http://mi.mathnet.ru/vsgtu1527}

%\crossref{10.14498/vsgtu1527}

%\elib{http://elibrary.ru/item.asp?id=29245104}



\Bibitem{Burm_Pros_IOP}

\by Burmasheva~N.~V.,  Prosviryakov~E.~Yu.

\paper Exact solutions for layered large-scale convection induced by tangential stresses specified on the free boundary of a fluid layer

\jour IOP Conference Series: Materials Science and Engineering

\yr 2017

\vol 208

\papernumber 012010

\crossref{10.1088/1757-899X/208/1/012010} 

 



\RBibitem{Burm_Pros_DREAM}

\by Бурмашева~Н.~В., Просвиряков~Е.~Ю.

\paper Точные решения для естественной конвекции слоистых течений вязкой несжимаемой жидкости при задании тангенциальных сил и линейного распределения температуры на границах слоя

\jour Diagnostics, Resource and Mechanics of Materials and Structures

\yr 2017

\issue 4

\pages 16--31

\crossref{10.17804/2410-9908.2017.4.016-031}

%\elib{http://elibrary.ru/item.asp?id=29245104}

%\misctransl

%\lang In Russian

%\by Burmasheva~N.~V.,  Prosviryakov~E.~Yu.

%\paper Exact solutions for natural convection of layered flows of a viscous incompressible fluid with specified tangential

%forces and the linear distribution of temperature on the layer boundaries

%\jour DREAM

%\yr 2017

%%\vol 21

%\issue 4

%\pages 16--31

%%\mathnet{http://mi.mathnet.ru/vsgtu1527}

%\crossref{10.17804/2410-9908.2017.4.016-031}

%%\elib{http://elibrary.ru/item.asp?id=29245104}



\Bibitem{Burm_Pros_AIP}

\by Burmasheva~N.~V.,  Prosviryakov~E.~Yu.

\paper Exact solution for the layered convection of a viscous incompressible fluid at specified temperature gradients and tangential forces on the free boundary

\jour AIP Conference Proceedings

\yr 2017

\vol 1915

\papernumber 040005

\crossref{10.1063/1.5017353}







\RBibitem{Gorsh_Pros1}

\by Горшков~А.~В., Просвиряков~Е.~Ю.

\paper Слоистая конвекция Бенара--Марангони при теплообмене по закону Ньютона--Рихмана

\jour Компьютерные исследования и моделирование

\yr 2016

\vol 8

\issue 6

\pages 927--940

\mathnet{http://mi.mathnet.ru/crm37}

%\misctransl

%\lang In Russian

%\by Gorshkov~A.~V.,  Prosviryakov~E.~Yu.

%\paper Layered Benard-Marangoni convection during heat transfer according to the Mewton's law of cooling

%\jour  Komp'yuternie issledovaniya i modelirovanie \rm[Computer Research and Modeling]

%\yr 2016

%\vol 8

%\issue 6

%\pages 927--940

%\mathnet{http://mi.mathnet.ru/crm37}



\RBibitem{Gorsh_Pros2}

\by Горшков~А.~В., Просвиряков~Е.~Ю.

\paper Аналитические решения стационарной сложной конвекции, описывающие поле касательных напряжений разного знака

\serial Тр. ИММ УрО РАН

\yr 2017

\vol 23

\issue 2

\pages 32--41

\mathnet{http://mi.mathnet.ru/timm1410}

\crossref{10.21538/0134-4889-2017-23-2-32-41}

\elib{http://elibrary.ru/item.asp?id=29295248}

%\misctransl

%\lang In Russian

%\by Gorshkov~A.~V.,  Prosviryakov~E.~Yu.

%\paper Analytic solutions of stationary complex convection describing a shear stress field of different signs

%\jour   Tr. IMM UrO RAN% \rm[???]

%\yr 2017

%\vol 23

%\issue 2

%\pages 32--41

%\crossref{10.21538/0134-4889-2017-23-2-32-41}



\RBibitem{Knyaz1}

\by Князев~Д.~В.

\paper Плоские течения вязкой бинарной жидкости между подвижными твердыми границами

\jour ПМТФ

\yr 2011

\vol 52

\issue 2

\pages 66--72

%\misctransl

%%\lang In Russian

%\by Knyazev~D.~V.

%\paper Two-dimensional flows of a viscous binary fluid between moving solid boundaries

%\jour J. Appl. Mech. Tech.Phys

%\yr 2011

%\vol 52

%\issue 2

%\pages 212--217

%\crossref{10.1134/S0021894411020088}







\RBibitem{Sidorov1}

\by Сидоров~А.~Ф.

\paper О двух классах решений уравнений механики жидкости и газа и их связи с теорией бегущих волн

\jour ПМТФ

\yr 1989

\vol 30

\issue 2

\pages 34--40

%\misctransl

%https://link.springer.com/article/10.1007%2FBF00852164

%\lang In Russian

%\by Sidorov~A.~F.

%\paper On two classes of solutions of the equations of fluid and gas mechanics and their relation to the theory of traveling waves

%\jour Prikladnaya mehanika i thenicheskaya fizika \rm [J. Appl. Mech. Tech.Phys]

%\yr 1989

%\vol 30

%\issue 2

%\pages 34--40



%\RBibitem{Sidorov2}

%\by Сидоров~А.~Ф.

%\paper Об одном классе решений уравнений газовой динамики и естественной конвекции

%\inbook Численные и аналитические методы решения задач механики сплошной среды

%\yr 1988

%\publaddr Свердловск

%\publ УНЦ АН СССР

%\yr 1981

%\eds А.~Ф.~Сидоров, Ю.~Н.~Кондюрин

%\pages 101--117

%\zmath{https://zbmath.org/?q=an:0512.76072}

%\mathscinet{https://mathscinet.ams.org/mathscinet-getitem?mr=0664611}



\RBibitem{Schwartz1}

\by  Шварц~К.~Г.

\paper Плоскопараллельное адвективное течение в горизонтальном слое несжимаемой жидкости с твердыми границами

\jour Изв. РАН. МЖГ

\yr 2014

\issue 4

\pages 26--30

\elib{https://elibrary.ru/item.asp?id=21811477}

%\misctransl

%\lang In Russian

%\by Shvarts~K.~G.

%\paper Plane-parallel advective flow in a horizontal layer of an incompressible fluid with solid boundaries

%\jour Izvestiya Rossiyskoi akademii nauk.  Mehanika zhidkosti i gaza \rm [Izvestiya of Russian academy of sciences. Mechanics of fluid and gas]

%\yr 2014

%%\vol 2

%\issue 4

%\pages 26--30



\RBibitem{Arist2006}

\by Аристов~С.~Н., Шварц~К.~Г.

\book Вихревые течения адвективной природы во вращающемся слое жидкости

\yr 2006

\publ Перм. гос. ун-т

\publaddr Пермь

\totalpages 155

%\misctransl

%\by Aristov~S.~N., Shvarts~K.~G.

%\book Vikhrevye techeniia advektivnoi prirody vo vrashchaiushchemsia

%sloe zhidkosti \rm [Vortical Flows of Advective Nature in a Rotating Fluid Layer]

%\yr 2006

%\publ Perm

%State Univ.

%\publaddr Perm

%\lang In Russian

%\totalpages 155



\RBibitem{Aristov_Schwartz4}

\by Аристов~С.~Н., Шварц~К.~Г.

\paper Адвективное течение во вращающейся жидкой пленке

\jour ПМТФ

\yr 2016

\vol 57

\issue 1

\pages 216--223

\crossref{10.15372/PMTF20160121}

%\misctransl

%%\lang In Russian

%\by Aristov~S.~N., Shvarts~K.~G.

%\paper Advective flow in a rotating liquid film

%\jour J. Appl. Mech. Tech.Phys

%\yr 2016

%\vol 57

%\issue 1

%\pages 188--194

%\crossref{10.1134/S0021894416010211}



\RBibitem{Schwartz2}

\by  Шварц~К.~Г.

\paper Влияние вращения на устойчивость адвективного течения в горизонтальном слое жидкости при малом значении числа Прандтля

\jour Изв. РАН. МЖГ

\yr 2005

\issue 2

\pages 29--38

%\misctransl

%\lang In Russian

%\by Shvarts~K.~G.

%\paper The effect of rotation on the stability of advective flow in a horizontal fluid layer with a small Prandtl number

%\jour Izvestiya Rossiyskoi akademii nauk.  Mehanika zhidkosti i gaza \rm [Izvestiya of Russian academy of sciences. Mechanics of fluid and gas]

%\yr 2005

%%\vol 2

%\issue 2

%\pages 29--38



\RBibitem{Schwartz3}

\by  Шварц~К.~Г.

\paper Моделирование крупномасштабных и мезомасштабных процессов в бароклинной атмосфере и океане

\jour Географический вестник

\yr 2013

\issue 1(24)

\pages 72--77

\elib{https://elibrary.ru/item.asp?id=19395931}

%%\zmath{https://?????}

%\misctransl

%\lang In Russian

%\by Shvarts~K.~G.

%\paper Modeling of large-scale and meso-scale processes in the baroclinic atmosphere and ocean

%\jour Geographicheskiy vestnik \rm [Geographical Bulletin]

%\yr 2013

%%\vol 2

%\issue 1(24)

%\pages 72--77





\Bibitem{Shtern}

\by Shtern~V.

\book Counterflows. Paradoxical Fluid Mechanics Phenomena

\publ Cambridge University Press

\publaddr  Cambridge 

\yr 2012

\totalpages 469~pp \nofrills

\crossref{10.1017/CBO9781139226516}



\Bibitem{Dekart}

\by Descartes~R.

\book La g\'eom\'etrie. Nouvelle \'ed. Avec le portrait de Descartes d'apr\`es Frans Hals

\zmath{https://zbmath.org/?q=an:02580439}

\publaddr Paris

\publ J.~Hermann

\totalpages 91

\yr 1927

\lang In French







\end{thebibliography}





\newpage

\makeentitle



\vspace{-5mm}



\AdditionalContent{Competing interests}{We declare that we have no conflicts of interests in the authorship and publication of this article.}

\AdditionalContent{Authors' contributions and responsibilities}{We are fully responsible for submitting the final manuscript in print. Each of us has approved the final version of manuscript.}

\AdditionalContent{Funding}{This work was supported by the Foundation for Assistance to Small Innovative

Enterprises in Science and Technology (the UMNIK program), agreement no.~12281GU/2017.}









\begin{thebibliography}{99} 



\Bibitem{Joseph:n}

\by Joseph~D.~D.

\book Stability of fluid motions

\yr 1976

\publ Springer–Verlag

\publaddr Berlin, Heidelberg, New York

\totalpages 282





\Bibitem{Gershuni1989:en}

\lang In Russian

\by Gershuni~G.~Z., Zhukhovitskii~E.~M., Nepomniashchii~A.~A.

\book Ustoichivost' konvektivnykh techenii

\yr 1989

\publ Nauka

\publaddr Moscow

\totalpages 320

%http://adsabs.harvard.edu/abs/1989MoIzN....R....G





\Bibitem{Pukhnachev4:en}

\paper Hierarchy of Models in Convection Theory

\by Pukhnachov~V.~V.

\jour J. Math. Sci. (N. Y.)

\yr 2004

\vol 123

\issue 6

\pages 4607--4620

\crossref{10.1023/B:JOTH.0000041478.45024.64}



\Bibitem{Ostroumov:en}

\lang In Russian

\by Ostroumov~G.~A.

\book Svobodnaia konvektsiia v usloviiakh vnutrennei zadachi \rm [Free Convection Under the Conditions of the Internal Problem]

\yr 1952

\publ Gostekhizdat

\publaddr Moscow

\totalpages 256





\Bibitem{Birikh:en}

\by Birikh~R.~V.

\paper Thermocapillary convection in a horizontal layer of liquid

\jour J. Appl. Mech. Tech. Phys.

\crossref{10.1007/BF00914697}

\vol 7

\issue 3

\pages 43–44 





\Bibitem{Smith:en}

\by Smith~M.~K., Davis~S.~H.

\paper Instabilities of dynamic thermocapillary liquid layers. Pt. 1. Convective instabilities

\jour J. Fluid Mech.

\yr 1983

\vol 132

\pages 119--144

\crossref{10.1017/S0022112083001512}



\Bibitem{Ortiz:en}

\by Ortiz-P\'{e}rez~A.~S., D\'{a}valos-Orozco~L.~A.

\paper Convection in a horizontal fluid layer under an inclined temperature gradient

\jour Phys. Fluids

\yr 2011

\vol 23

\issue 8

\pages 084107–084111

\crossref{10.1063/1.3626009}





\Bibitem{Andreev:en}

\lang In Russian

\by Andreev~V.~K.

\book Resheniia Birikha uravnenii konvektsii i nekotorye ego obobshcheniia \rm [Birikh Solutions to Convection Equations and Some of its Extensions]

\yr 2010

\publ IVM SO RAN Publ.

\publaddr Krasnoyarsk

\totalpages 68



\Bibitem{Andreev_bek:en}

\by Andreev~V.~K., Bekezhanova~V.~B.

\paper Stability of non-isothermal fluids (Review)

\jour J. Appl. Mech. Tech. Phys.

\yr 2013

\vol 54

\issue 2

\pages 171--184

\crossref{10.1134/S0021894413020016}





\Bibitem{Andreev_step:en}

\paper Unidirectional flows of binary mixtures within the framework of the Oberbeck–Boussinesq model

\by Andreev~V.~K., Stepanova~I.~V.

\jour Fluid Dyn.

\yr 2016

\vol 51

\issue 2

\pages 136--147

\elib{https://elibrary.ru/item.asp?id=27081601}

\crossref{10.1134/S0015462816020022}



\Bibitem{Bek_Gon_Rez_Shef:en}

\paper Stability of two-layer fluid flows with evaporation at the interface

\by Bekezhanova~V.~B., Shefer~I.~A., Goncharova~O.~N., Rezanova~E.~B.

\jour Fluid Dyn.

\yr 2017

\vol 52

\issue 2

\pages 189--200

\crossref{10.1134/S001546281702003X}

\elib{https://elibrary.ru/item.asp?id=29492139}



\Bibitem{Bir_Den_Kost1:en}

\paper The development of Marangoni convection induced by local addition of a surfactant

\by Birikh~R.~V., Denisova~M.~O., Kostarev~K.~G.

\jour Fluid Dyn.

\yr 2011

\vol 46

\issue 6

\pages 890--900

\elib{https://elibrary.ru/item.asp?id=18330572}

\crossref{10.1134/S0015462811060068}





\Bibitem{Bir_Den_Kost2:en}

\by Birih~R.~V., Denisova~M.~O., Kostarev~K.G.

\paper Development of concentration-capillary convection on an interfacial surface

\jour 	Fluid Dyn.

\yr 2015

\vol 50

\issue 3

\pages 361--370

\crossref{10.1134/S0015462815030060}

\elib{https://elibrary.ru/item.asp?id=23984641}



\Bibitem{Bir_Pukhn:en}

\by Birih~R.~V., Puhnachev~V.~V.

\paper An axial convective flow in a rotating tube with a longitudinal temperature gradient

\jour Dokl. Phys.

\yr 2011

\vol 56

\issue 1

\pages 47--52

\crossref{10.1134/S1028335811010095}



\Bibitem{Birikh2:en}

\by Birih~R.~V., Puhnachev~V.~V.

\paper Convective flow in a horizontal channel with non-newtonian surface rheology under time-dependent longitudinal temperature gradient

\jour Fluid Dyn.

\yr 2015

\vol 50

\issue 1

\pages 173--179

\crossref{10.1134/S0015462815010172}

\elib{https://elibrary.ru/item.asp?id=24010631}





\Bibitem{Gon_Kab2:en}

\by Goncharova~O.N., Kabov~O.A.

\paper Gravitational-thermocapillary convection of fluid in the horizontal layer in co-current gas flow

\jour Dokl. Phys.

\yr 2009

\vol 54

\issue 5

\pages 242--247

\crossref{10.1134/S1028335809050061}



\Bibitem{Gon_rez:en}

\by Goncharova~O.~N., Rezanova~E.~V.

\paper 	Example of an exact solution of the stationary problem of two-layer flows with evaporation at the interface

\jour J. Appl. Mech. Tech. Phys.

\yr 2014

\vol 55

\issue 2

\pages 247--257

\crossref{10.1134/S0021894414020072}





\Bibitem{Pukhnachev6:rm}

\lang In Russian

\by Puhnachev~V.~V.

\paper Non-stationary analogues of the Birikh solution

\jour  Izvestiya AltGU 

\yr 2011

\issue 1--2

\pages 62--69





\Bibitem{Aristov_avto:en}

\lang In Russian

\by Aristov~S.~N.

\thesis Vortex flows in thin layers of a liquid

\thesisinfo Dr. Sci. Thesis (Phys. \& Math.)

\yr 1990

\publaddr Vladivostok

\totalpages 32





\Bibitem{Aristov_Pros3:en}

\by Aristov~S.N., Prosviryakov~E.~Yu.

\paper Nonuniform convective Couette flow

\jour Fluid Dyn.

\yr 2016

\vol 51

\issue 5

\pages 581--587

\crossref{10.1134/S001546281605001X}



\Bibitem{Aristov_Pros5:en}

\by Aristov~S.N., Prosviryakov~E.~Yu.

\paper Unsteady layered vortical fluid flows

\jour Fluid Dyn.

\yr 2016

\vol 51

\issue 2

\pages 148--154

\crossref{10.1134/S0015462816020034}

\elib{https://elibrary.ru/item.asp?id=27097904}





\Bibitem{Aristov_Pros7:en}

\lang In Russian

\by Aristov~S.~N., Prosviryakov~E.~Yu.

\paper On laminar flows of planar free convection

\jour Nelin. Dinam.

\yr 2013

\vol 9

\issue 4

\pages 651--657

\crossref{10.20537/nd1304004}





\Bibitem{Aristov_Pros_Spev1:en}

\by Aristov~S.~N., Prosviryakov~E.~Yu., Spevak~L.F.

\paper Unsteady-state B\'enard–Marangoni convection in layered viscous incompressible flows

\jour  Theor. Found. Chem. Eng.

\yr 2016

\vol 50

\issue 2

\pages 137--146

\crossref{10.1134/S0040579516020019}

\elib{https://elibrary.ru/item.asp?id=27053288}







\Bibitem{Aristov_Pros_Spev2:en}

\lang In Russian

\crossref{10.7242/1999-6691/2015.8.4.38}

\by Aristov~S.~N., Prosviryakov~E.~Yu., Spevak~L.F.

\paper Nonstationary laminar thermal and solutal Marangoni convection of a viscous fluid

\jour  Computational Continuum Mechanics

\yr 2015

\vol 8

\issue 4

\pages 445--456









\RBibitem{Burm_Pros_Samara:en}

\lang In Russian

\by Burmasheva~N.~V., Prosviryakov~E.~Yu.

\paper A large-scale layered stationary convection of~a~incompressible viscous fluid under the action

of shear stresses at the upper boundary. Velocity field investigation

\jour Vestn. Samar. Gos. Tekhn. Univ., Ser. Fiz.-Mat. Nauki \rm [J. Samara State Tech. Univ., Ser. Phys. Math. Sci.]

\yr 2017

\vol 21

\issue 1

\pages 180--196

\crossref{10.14498/vsgtu1527}





\Bibitem{Burm_Pros_IOP:en}

\by Burmasheva~N.~V.,  Prosviryakov~E.~Yu.

\paper Exact solutions for layered large-scale convection induced by tangential stresses specified on the free boundary of a fluid layer

\jour IOP Conference Series: Materials Science and Engineering

\yr 2017

\vol 208

\papernumber 012010

\crossref{10.1088/1757-899X/208/1/012010} 

 



\Bibitem{Burm_Pros_DREAM:en}

\lang In Russian

\by Burmasheva~N.~V.,  Prosviryakov~E.~Yu.

\paper Exact solutions for natural convection of layered flows of a viscous incompressible fluid with specified tangential forces and the linear distribution of temperature on the layer boundaries

\jour Diagnostics, Resource and Mechanics of Materials and Structures

\yr 2017

\issue 4

\pages 16--31

\crossref{10.17804/2410-9908.2017.4.016-031}



\Bibitem{Burm_Pros_AIP:en}

\by Burmasheva~N.~V.,  Prosviryakov~E.~Yu.

\paper Exact solution for the layered convection of a viscous incompressible fluid at specified temperature gradients and tangential forces on the free boundary

\jour AIP Conference Proceedings

\yr 2017

\vol 1915

\papernumber 040005

\crossref{10.1063/1.5017353}







\Bibitem{Gorsh_Pros1:en}

\lang In Russian

\by Gorshkov~A.~V., Prosviryakov~E.~Yu.

\paper Layered B\'enard-Marangoni convection during heat transfer according to the Newton's law of cooling

\jour Computer Research and Modeling

\yr 2016

\vol 8

\issue 6

\pages 927--940



\Bibitem{Gorsh_Pros2:en}

\lang In Russian

\by Gorshkov~A.~V., Prosviryakov~E.~Yu.

\paper Analytic solutions of stationary complex convection describing a shear stress field of different signs

\serial Trudy Inst. Mat. i Mekh. UrO RAN

\yr 2017

\vol 23

\issue 2

\pages 32--41

\crossref{10.21538/0134-4889-2017-23-2-32-41}



\Bibitem{Knyaz1:en}

\by Knyazev~D.~V.

\paper Two-dimensional flows of a viscous binary fluid between moving solid boundaries

\jour J. Appl. Mech. Tech. Phys.

\yr 2011

\vol 52

\issue 2

\pages 212--217

\crossref{10.1134/S0021894411020088}







\Bibitem{Sidorov1:en}

\jour J. Appl. Mech. Tech. Phys. 

\yr 1989

\crossref{10.1007/BF00852164}

\vol 30

\issue 2

\pages 197–203 

\paper Two classes of solutions of the fluid and gas mechanics equations and their connection to traveling wave theory

\by Sidorov~A.~F.





%\Bibitem{Sidorov2:en}

%\lang In Russian

%\by Sidorov~A.~F.

%\paper On a class of solutions of equations of gas dynamics and natural convection

%\inbook Chislennye i analiticheskie metody resheniia zadach mekhaniki sploshnoi sredy \rm [Numerical and analytical methods for solving problems of continuum mechanics]

%\eds A.~F.~Sidorov, Yu.~N.~Kondyurin

%\yr 1988

%\publaddr Sverdlovsk

%\publ UNTs AN SSSR Publ.

%\yr 1981

%\pages 101--117







\Bibitem{Schwartz1:en}

\by Schwarz~E.~G. 

\crossref{10.1134/S0015462814040036}

\jour Fluid Dyn.

\yr 2014

\vol 49

\issue 4

\pages 438–442 

\paper Plane-parallel advective flow in a horizontal incompressible fluid layer with rigid boundaries





\Bibitem{Arist2006:en}

\lang In Russian

\by Aristov~S.~N., Shvarts~K.~G.

\book Vikhrevye techeniia advektivnoi prirody vo vrashchaiushchemsia sloe zhidkosti \rm [Vortical Flows of Advective Nature in a Rotating Fluid Layer]

\yr 2006

\publ Perm State Univ.

\publaddr Perm

\totalpages 155





\Bibitem{Aristov_Schwartz4:en}

\by Aristov~S.~N., Shvarts~K.~G.

\paper Advective flow in a rotating liquid film

\jour J. Appl. Mech. Tech.Phys

\yr 2016

\vol 57

\issue 1

\pages 188--194

\crossref{10.1134/S0021894416010211}



\Bibitem{Schwartz2:en}

\crossref{10.1007/s10697-005-0059-7}

\jour Fluid Dyn.

\yr 2005

\vol 40

\issue 2

\pages 193–201 

\paper Effect of rotation on the stability of advective flow in a horizontal fluid Layer at a small Prandtl number

\by Schwarz~K.~G.





\Bibitem{Schwartz3:en}

\lang In Russian

\by Shvarts~K.~G.

\paper Modeling of large-scale and meso-scale processes in the baroclinic atmosphere and ocean

\jour Geographical Bulletin

\yr 2013

\issue 1(24)

\pages 72--77





\Bibitem{Shtern:en}

\by Shtern~V.

\book Counterflows. Paradoxical Fluid Mechanics Phenomena

\publ Cambridge University Press

\publaddr  Cambridge 

\yr 2012

\totalpages 469~pp \nofrills

\crossref{10.1017/CBO9781139226516}



\Bibitem{Dekart:en}

\by Descartes~R.

\book La g\'eom\'etrie. Nouvelle \'ed. Avec le portrait de Descartes d'apr\`es Frans Hals

\zmath{https://zbmath.org/?q=an:02580439}

\publaddr Paris

\publ J.~Hermann

\totalpages 91

\yr 1927

\lang In French







\end{thebibliography}



\renewcommand{\baselinestretch}{0.9}





%\end{document}

